\documentclass{amsart}

\usepackage[T1]{fontenc}
\usepackage{amsmath,amssymb,mathtools}
\usepackage[margin=0.78in]{geometry}
\usepackage{microtype}
\usepackage[hidelinks]{hyperref}

\newtheorem{theorem}{Theorem}[section]
\newtheorem{lemma}[theorem]{Lemma}

\newcommand{\F}{\mathbb{F}}
\newcommand{\Gammavec}{\Gamma}
\newcommand{\ind}{\mathbf{1}}
\newcommand{\im}{\operatorname{im}}

\title{A Proof of the Cycle Double Cover Conjecture}
\author{OPENAI}
\date{}

\begin{document}

\begin{abstract}
We prove the cycle double cover conjecture, posed by Tutte, Itai and Rodeh,
Szekeres, and Seymour, which asserts that every bridgeless undirected graph has
a collection of cycles which covers every edge exactly twice.
\end{abstract}

\maketitle

\section{Introduction}

A \emph{cycle double cover} of a graph is a multiset of cycles in which every
edge occurs exactly twice. The conjecture was posed by Tutte~\cite{Tutte1987},
Itai and Rodeh~\cite{ItaiRodeh1978}, Szekeres~\cite{Szekeres1973}, and
Seymour~\cite{Seymour1979}. See the survey of Jaeger~\cite{Jaeger1985} for
further background.

\begin{theorem}
Every finite bridgeless undirected graph has a cycle double cover.
\end{theorem}

Among partial results, Jaeger observed that the conjecture holds for planar
graphs by taking the facial boundary cycles of their blocks
\cite[Sections~2.1 and~3.1]{Jaeger1985}, Szekeres observed that it holds for
3-edge-colourable cubic graphs by taking the three unions of pairs of colour
classes \cite[p.~367]{Szekeres1973}, and Alspach, Goddyn, and Zhang proved it
for bridgeless graphs with no Petersen subdivision~\cite{AlspachGoddynZhang1994}.
The proof proceeds by first using that it suffices (as is standard) to consider
cubic graphs. Then, using the 8-flow theorem and a result of Tutte, we have a
labeling of the edges by nonzero elements of $\Gamma=\F_2^3$ such that the sum
at each vertex is zero. The key reduction is then to convert this labeling into
a labeling of the edges by sets of two elements in $\Gamma$ such that each
element of $\Gamma$ appears either zero or twice next to a given vertex. This
reduction eventually reduces to an elementary linear algebra argument.

\medskip
\noindent\textbf{Statement of AI use.}
The proof in this note is entirely due to GPT 5.6 Sol Ultra and the writeup
with Codex (with GPT 5.6 Sol).

\section{Proof of the conjecture}

We allow parallel edges and regard two parallel edges as a cycle. By
Jaeger~\cite[Proposition~4]{Jaeger1985}, it suffices to treat loopless cubic
graphs. In fact, Jaeger observed that a minimum counterexample must fail to be
3-edge-colourable---that is, it must be a snark---and we return to this
observation below.

Fix an orientation of a graph. If $A$ is an abelian group, an $A$-flow is a
map $f:E(G)\to A$ for which, at every vertex, the sum on edges directed out
equals the sum on edges directed in. It is \emph{nowhere-zero} if $f(e)\ne 0$
for every edge. For an integer $k\ge 2$, an \emph{integer $k$-flow} is an
integer-valued flow $\phi$ satisfying $0<|\phi(e)|<k$ on every edge.

Let $\Gamma=\F_2^3$, written additively. Kilpatrick and Jaeger independently
proved that every bridgeless graph has a nowhere-zero $\Gamma$-flow
\cite{Kilpatrick1975,Jaeger1976}, or, equivalently by Tutte's group-flow
theorem, a nowhere-zero 8-flow~\cite{Tutte1954}. (Seymour's later 6-flow
theorem~\cite{Seymour1981} is stronger, but unnecessary here.) We now massage
this $\Gamma$-flow into a cycle double cover; this reduction only requires that
the underlying graph $G$ is loopless and cubic.

\begin{lemma}\label{lem:pair-labeling}
Let $G$ be a loopless cubic multigraph. Suppose that every edge $e$ is assigned
a two-element set $P_e\subseteq\Gamma$ such that, for every $v\in V(G)$ and
$s\in\Gamma$,
\begin{equation}\label{eq:local-evenness}
  \bigl|\{e\ni v:s\in P_e\}\bigr|\in\{0,2\}.
\end{equation}
Then $G$ has a cycle double cover.
\end{lemma}

A proper 3-edge-colouring is an assignment of this form. For distinct
$e_1,e_2\in\Gamma$, assign $\{0,e_1\}$, $\{0,e_2\}$, and $\{e_1,e_2\}$ to the
red, blue, and green edges, respectively. At each vertex each of $0,e_1,e_2$
then appears twice. Thus Lemma~\ref{lem:pair-labeling} can be seen as a suitably
relaxed variant of a proper 3-edge-colouring.

\begin{proof}
For $s\in\Gamma$, let $M_s=\{e:s\in P_e\}$. By
\eqref{eq:local-evenness}, every vertex has degree zero or two in $M_s$, so
$M_s$ is a disjoint union of cycles. Every edge belongs to exactly two of the
$M_s$, since $P_e$ has two elements. The cycle components of all the $M_s$,
taken with multiplicity, form a cycle double cover.
\end{proof}

It remains to construct the sets $P_e$. Fix a nowhere-zero $\Gamma$-flow $f$.
At each vertex $v$, locally order the incident edges as $a,b,c$ and write
$x=f(a)$, $y=f(b)$, and $z=f(c)$. Since $\Gamma$ has characteristic two, the
flow equation is $x+y+z=0$, so $z=x+y$; moreover $x$ and $y$ are distinct.
Define
\begin{equation}\label{eq:g-definition}
  g_{v,a}=0,\qquad g_{v,b}=x,\qquad g_{v,c}=0.
\end{equation}
For any $t\in\Gamma$, the three local sets
\begin{equation}\label{eq:local-sets}
  \{t+g_{v,e},\,t+g_{v,e}+f(e)\}\qquad(e\ni v)
\end{equation}
are $\{t,t+x\}$, $\{t+x,t+z\}$, and $\{t,t+z\}$. Hence every vector occurs in
zero or two of them. This assignment works locally, but the two ends of an
edge need not assign it the same set.

For $e=uv$, put $d_e=g_{u,e}+g_{v,e}$. For any $p\in\Gamma$,
$\{A,A+p\}=\{B,B+p\}$ precisely when $A+B\in\{0,p\}$. Apply this with
$A=t_u+g_{u,e}$, $B=t_v+g_{v,e}$, and $p=f(e)$. Then
$A+B=t_u+t_v+d_e$, and the choice of $\epsilon_e\in\F_2$ records whether this
vector is $0$ or $f(e)$. Thus the local sets in \eqref{eq:local-sets} agree
across every edge precisely when there are $t_v\in\Gamma$ and
$\epsilon_e\in\F_2$ satisfying
\begin{equation}\label{eq:gluing-system}
  t_u+t_v+\epsilon_e f(e)=d_e\qquad(e=uv).
\end{equation}

\begin{lemma}\label{lem:linear-system}
The system \eqref{eq:gluing-system} has a solution.
\end{lemma}

Assuming the lemma, define
\[
  P_e=\{t_v+g_{v,e},\,t_v+g_{v,e}+f(e)\}
\]
using either endpoint $v$ of $e$. Equation~\eqref{eq:gluing-system} makes this
independent of the endpoint, and $f(e)\ne 0$ makes the two elements distinct.
The local calculation above gives \eqref{eq:local-evenness};
Lemma~\ref{lem:pair-labeling} proves the theorem.

\begin{proof}[Proof of Lemma~\ref{lem:linear-system}]
Let
\[
  L:\Gamma^{V(G)}\oplus\F_2^{E(G)}\longrightarrow\Gamma^{E(G)},
  \qquad
  L(t,\epsilon)_e=t_u+t_v+\epsilon_e f(e)\quad(e=uv).
\]
Thus \eqref{eq:gluing-system} asks whether $d=(d_e)_e$ belongs to $\im L$.
Let $\Gamma^*$ be the dual vector space of $\Gamma$; thus an element of
$\Gamma^*$ is an $\F_2$-linear map from $\Gamma$ to $\F_2$. We may write a
dual vector to $\Gamma^{E(G)}$ as a family
$\eta=(\eta_e)_{e\in E(G)}$ with $\eta_e\in\Gamma^*$. By taking duals,
$d\in\im L$ if and only if every such family which takes the value zero on
$\im L$ also satisfies $\sum_e\eta_e(d_e)=0$. Now
\[
  \sum_e\eta_e\bigl(L(t,\epsilon)_e\bigr)
  =\sum_v\left(\sum_{e\ni v}\eta_e\right)(t_v)
   +\sum_e\epsilon_e\eta_e(f(e)).
\]
Since the $t_v$ and $\epsilon_e$ may be chosen independently, this vanishes for
every $(t,\epsilon)$ precisely when
\begin{equation}\label{eq:dual-constraints}
  \eta_e(f(e))=0\quad(e\in E(G)),
  \qquad
  \sum_{e\ni v}\eta_e=0\quad(v\in V(G)).
\end{equation}
Consequently, it suffices to prove that every family satisfying
\eqref{eq:dual-constraints} also satisfies
\begin{equation}\label{eq:d-annihilated}
  \sum_e\eta_e(d_e)=0.
\end{equation}

Fix $v$ and retain the notation $a,b,c,x,y,z$. The conditions
\eqref{eq:dual-constraints} become
\begin{equation}\label{eq:local-dual-constraints}
  \eta_a+\eta_b+\eta_c=0,
  \qquad
  \eta_a(x)=0,
  \qquad
  \eta_b(y)=0,
  \qquad
  \eta_c(z)=0.
\end{equation}
Set $\lambda=\eta_b(x)$. Since $\eta_c=\eta_a+\eta_b$ and $z=x+y$,
\[
  0=\eta_c(z)=\eta_a(y)+\eta_b(x),
\]
where we used $\eta_a(x)=\eta_b(y)=0$. Hence $\eta_a(y)=\lambda$. By
\eqref{eq:g-definition}, only the edge $b$ contributes at $v$, and therefore
\begin{equation}\label{eq:local-g-sum}
  \sum_{e\ni v}\eta_e(g_{v,e})=\eta_b(x)=\lambda.
\end{equation}

We next interpret $\lambda$. If $\lambda=0$, all three dual vectors vanish on
$W=\langle x,y\rangle$. There is a unique nonzero dual vector which vanishes on
this two-dimensional space, so each of $\eta_a,\eta_b,\eta_c$ is either zero or
this vector. Since their sum is zero, the nonzero vector occurs zero or two
times. If $\lambda=1$, their values on the ordered basis $x,y$ are $(0,1)$,
$(1,0)$, and $(1,1)$, so all three are nonzero. In either case, $\lambda$ is
the parity of the number of nonzero members of
$\{\eta_a,\eta_b,\eta_c\}$. Thus, writing $\ind_{\eta\ne 0}$ for the
corresponding bit,
\begin{equation}\label{eq:local-parity}
  \sum_{e\ni v}\eta_e(g_{v,e})
  =\sum_{e\ni v}\ind_{\eta_e\ne 0}.
\end{equation}

Finally, for $e=uv$, the definition $d_e=g_{u,e}+g_{v,e}$ and linearity of
$\eta_e$ give
$\eta_e(d_e)=\eta_e(g_{u,e})+\eta_e(g_{v,e})$. Summing over all edges and then
grouping the two endpoint terms at each vertex, \eqref{eq:local-parity} gives
\[
  \sum_e\eta_e(d_e)
  =\sum_v\sum_{e\ni v}\eta_e(g_{v,e})
  =\sum_v\sum_{e\ni v}\ind_{\eta_e\ne 0}.
\]
Each edge with $\eta_e\ne 0$ occurs twice in the last sum, once at each
endpoint. Hence it equals $2\sum_e\ind_{\eta_e\ne 0}=0$ in $\F_2$. This proves
\eqref{eq:d-annihilated}; by the duality criterion above, $d\in\im L$, so
\eqref{eq:gluing-system} has a solution.
\end{proof}

\begin{thebibliography}{10}

\bibitem{AlspachGoddynZhang1994}
B.~Alspach, L.~A. Goddyn, and C.-Q. Zhang,
\emph{Graphs with the circuit cover property},
Trans. Amer. Math. Soc. \textbf{344} (1994), 131--154.

\bibitem{ItaiRodeh1978}
A.~Itai and M.~Rodeh,
\emph{Covering a graph by circuits},
in \emph{Automata, Languages and Programming} (G.~Ausiello and C.~B\"ohm, eds.),
Lecture Notes in Comput. Sci. \textbf{62}, Springer, 1978, 289--299.

\bibitem{Jaeger1985}
F.~Jaeger,
\emph{A survey of the cycle double cover conjecture},
in \emph{Cycles in Graphs}, North-Holland Math. Stud. \textbf{115},
North-Holland, 1985, 1--12.

\bibitem{Jaeger1976}
F.~Jaeger,
\emph{On nowhere-zero flows in multigraphs},
in \emph{Proceedings of the Fifth British Combinatorial Conference},
Congr. Numer. XV, Utilitas Math., 1976, 373--378.

\bibitem{Kilpatrick1975}
P.~A. Kilpatrick,
\emph{Tutte's first colour-cycle conjecture},
M.Sc. thesis, University of Cape Town, 1975.

\bibitem{Seymour1979}
P.~D. Seymour,
\emph{Sums of circuits},
in \emph{Graph Theory and Related Topics}, Academic Press, 1979, 341--355.

\bibitem{Seymour1981}
P.~D. Seymour,
\emph{Nowhere-zero 6-flows},
J. Combin. Theory Ser. B \textbf{30} (1981), 130--135.

\bibitem{Szekeres1973}
G.~Szekeres,
\emph{Polyhedral decompositions of cubic graphs},
Bull. Austral. Math. Soc. \textbf{8} (1973), 367--387.

\bibitem{Tutte1954}
W.~T. Tutte,
\emph{A contribution to the theory of chromatic polynomials},
Canad. J. Math. \textbf{6} (1954), 80--91.

\bibitem{Tutte1987}
W.~T. Tutte,
personal correspondence with H.~Fleischner, July 22, 1987.

\end{thebibliography}

\end{document}
